\section{Integration with GFM Machinery}
\label{sec:integration}

This section connects the exogenous verification construction of
\S\ref{sec:gfm_claims} to the existing machinery of
Papers~1--4, showing how committed claims improve the fidelity of each
mechanism.

\subsection{Closing the Observation-Channel Integrity Assumption}
\label{sec:channel_independence}

Paper~4's multi-channel attribution (\S2.4 of~\citet{lasser2026scm})
aggregates evidence from multiple observation channels via log-odds
combination.  The aggregation's validity depends on the channels being
\emph{independently informative}: each channel's evidence is
conditionally independent of the others given the true state.  If an
agent can write to multiple channels simultaneously---because they all
run on the same substrate---the independence assumption fails, and the
log-odds aggregation can be arbitrarily manipulated.

The committed-claim construction closes this gap.  When each observation
channel's verification is rooted in a substrate-exclusive witness
(Definition~\ref{def:alg_witness}), and the witnesses are on different
substrates, a coalition on any single substrate cannot correlate multiple
channels.

\begin{proposition}[Channel Independence under Committed Verification]
\label{prop:channel_independence}
Let $\{E_1, \ldots, E_n\}$ be observation channels for agent $a$, and
let each $E_i$ be verified by an algorithmic witness
$\Witness^{(i)}_{\mathrm{alg}}$ on substrate $s_i$.  If the witnesses
are on distinct substrates ($s_i \neq s_j$ for $i \neq j$) and each
$s_i \neq s(a)$, then the verified outputs $\{v_1, \ldots, v_n\}$ are
conditionally independent given $a$'s true capability profile
$\mathbf{C}_a$.
\end{proposition}

\begin{proof}
Each $v_i = \pi_{B_i}(a, \Witness^{(i)}_{\mathrm{alg}})$ is determined
by $a$'s behavior under exercise protocol $\pi_{B_i}$ and the witness's
deterministic execution (Definition~\ref{def:alg_witness}).  Agent $a$
cannot influence $\Witness^{(i)}_{\mathrm{alg}}$'s computation
(substrate exclusivity).  Witness $\Witness^{(i)}_{\mathrm{alg}}$
cannot influence $\Witness^{(j)}_{\mathrm{alg}}$'s computation ($s_i
\neq s_j$).  Therefore, conditional on $a$'s true behavior
(determined by $\mathbf{C}_a$), the verdicts $v_i$ are functions of
independent computations on disjoint substrates, hence conditionally
independent.
\end{proof}

\begin{remark}
The proposition requires $n$ distinct witness substrates for $n$
channels.  When the substrate partition has $m$ substrates total, at
most $m - 1$ independent channels can be constructed (one substrate is
occupied by $a$, and each remaining substrate hosts at most one
independent witness).  This bounds the effective multiplicity of
independent verification channels by $m - 1$, connecting verification
depth to substrate diversity.
\end{remark}

\subsection{Committed Baselines for SPRT Detection}
\label{sec:committed_baselines}

Paper~4's SPRT-based scorpion detection compares an agent's behavioral
observations against a baseline.  In the current framework, the baseline
is the \emph{actor's own estimate} of the agent's expected behavior---an
endogenous quantity that is subject to the actor's model uncertainty and
potential manipulation.

Committed capability claims (\S\ref{sec:capability_claims}) provide an
exogenous alternative.  When agent $a_j$ has committed to capability
profile $\mathbf{C}_a$ and the commitment has been verified by a
substrate-exclusive witness, the behavioral distribution $p(\cdot \mid
\mathbf{C}_a)$ consistent with the committed profile becomes the
baseline.  The SPRT then tests whether $a_j$'s observed behavior is
consistent with its own committed claims, rather than with the actor's
model of what $a_j$ should be doing.

\begin{proposition}[Detection Improvement under Committed Baselines]
\label{prop:committed_detection}
Let $p_0$ be the behavioral distribution under the committed capability
profile $\mathbf{C}_a$, and let $\hat{p}_0$ be the actor's estimated
baseline.  If the agent's true distribution is $p_1 \neq p_0$, then
the SPRT detection rate satisfies:
\[
    \frac{1}{\mathbb{E}_{p_1}[T_{\mathrm{detect}}^{\mathrm{committed}}]}
    \;\geq\;
    \frac{1}{\mathbb{E}_{p_1}[T_{\mathrm{detect}}^{\mathrm{estimated}}]}
\]
whenever $D_{\mathrm{KL}}(p_1 \| p_0) \geq D_{\mathrm{KL}}(p_1 \|
\hat{p}_0)$, i.e., when the committed baseline is at least as
discriminating as the estimated baseline.  Equality holds if and only if
$\hat{p}_0 = p_0$.
\end{proposition}

\begin{proof}
By Proposition~\ref{prop:sprt_detection}, the expected detection time
is inversely proportional to the KL divergence between the true and
baseline distributions.  If $D_{\mathrm{KL}}(p_1 \| p_0) \geq
D_{\mathrm{KL}}(p_1 \| \hat{p}_0)$, then $\mathbb{E}[T^{\mathrm{committed}}]
\leq \mathbb{E}[T^{\mathrm{estimated}}]$, and detection is faster
(or equally fast) under the committed baseline.
\end{proof}

The committed baseline is at least as discriminating as the estimated
baseline whenever the agent's commitment is \emph{informative}---that
is, when the committed profile $\mathbf{C}_a$ constrains the expected
behavioral distribution more tightly than the actor's prior estimate
does.  Since agents commit to specific capability claims (not vague
ranges), committed baselines are generically more informative than
actor-estimated baselines, especially during the initial-period
noise that Paper~1's trust EWMA is most vulnerable to.

\subsection{Trust Dynamics with Verified Claims}
\label{sec:trust_dynamics}

Paper~1's trust EWMA $T_j$ updates from prediction
residuals~\citep{lasser2026gfm}.  Committed and verified claims anchor
the prediction baseline and affect trust dynamics in three ways.

\paragraph{Reduced initial-period noise.}
A new agent entering a population with no trust history begins with
$T_j$ at a prior value, and the EWMA requires multiple observations to
converge.  Verified capability claims provide a shortcut: a commitment
that is verified by a substrate-exclusive witness establishes a baseline
that the actor's trust model can condition on, reducing the number of
observations needed for convergence.  The verified claim does not
replace the behavioral evidence---it provides a structural prior that
the EWMA refines.

\paragraph{Falsification as accelerated detection.}
When a commitment is falsified ($v = \texttt{fail}$), the prediction
residual is large by construction: the agent predicted it would
demonstrate capability $c_k$ and did not.  In the EWMA framework, this
produces a correspondingly large negative trust update.  The key
structural property is that falsification is \emph{binary} (pass/fail
against an algorithmic protocol), not graded (how consistent is
behavior?), so the signal is clean and unambiguous.  A single falsified
commitment can move $T_j$ more than many observations of mildly
inconsistent behavior.

\paragraph{Exogenous risk residuals.}
Paper~4's risk trust $\Trisk_j$ is computed from the residual between
claimed and realized risk.  Currently, both the claim and the
realization are endogenous to the actor's model.  Committed risk claims
(\S\ref{sec:risk_claims}), verified by cross-substrate witnesses,
replace the endogenous claim with an exogenous one: the residual
$\mathbf{r}^{\mathrm{risk}}_j$ is now between what the agent committed
to and what the witness measured, not between what the actor estimated
and what the actor observed.  This sharpens $\Trisk_j$ by removing one
of the two endogenous quantities from the residual computation.

\subsection{The Anti-Monopolar Verification Theorem}
\label{sec:anti_monopolar}

Substrate-exclusive witnesses require the substrate partition to have
$m \geq 2$~(\S\ref{sec:witness_m2}).  The following theorem makes the
consequence precise.

\begin{theorem}[Anti-Monopolar Verification]
\label{thm:anti_monopolar}
Let $\mathcal{A} = \{a_1, \ldots, a_n\}$ be a population of agents
with substrate partition $\{S_1, \ldots, S_m\}$.  If $m = 1$ (the
population is monopolar), then for every agent $a_j \in \mathcal{A}$
and every claim $C$, the set of substrate-exclusive witnesses is empty:
\[
    \Witness^{\perp}(a_j, C)
    \;\triangleq\;
    \{w : s(w) \neq s(a_j)\}
    \;=\; \emptyset.
\]
Consequently, every evaluation function $E$ available to the population
satisfies $\rho(a_j, E) > 0$, and all verification is endogenous.
\end{theorem}

\begin{proof}
If $m = 1$, then all agents (including any candidate witness $w$)
occupy the same substrate: $s(w) = s(a_j)$ for all $w \in
\mathcal{A}$.  By Definition~\ref{def:alg_witness}, an algorithmic
witness requires $s(w) \neq s(a_j)$, so no member of $\mathcal{A}$
qualifies.  Since the population has no access to substrates outside
$S_1$ (by the monopolarity assumption), no process outside
$\mathcal{A}$ on a different substrate is available either.  Therefore
$\Witness^{\perp}(a_j, C) = \emptyset$ for all $a_j$ and $C$.

With no substrate-exclusive witnesses, every evaluator $E$ available to
the population runs on substrate $S_1$.  Agent $a_j$, also on $S_1$,
has $\mathcal{R}(a_j, E) \supseteq \{x \in \mathcal{I}(E) :
x \text{ is addressable from } S_1\}$, which is non-empty (the agent
and evaluator share a substrate and can therefore interact).  Thus
$\rho(a_j, E) > 0$ for all available $E$.
\end{proof}

\begin{corollary}[Monopolar Undetectability]
\label{cor:monopolar_undetect}
In a monopolar population ($m = 1$), no agent can distinguish between:
\begin{enumerate}
    \item[(a)] Verification passing because the claim is true, and
    \item[(b)] Verification passing because the evaluator is endogenous
    and the claimant has influenced the evaluation.
\end{enumerate}
\end{corollary}

\begin{proof}
Distinguishing between (a) and (b) requires a reference evaluation
that is exogenous to the claimant.  By Theorem~\ref{thm:anti_monopolar},
no such evaluation exists when $m = 1$.  Therefore the distinction is
unresolvable from within the population.
\end{proof}

Theorem~\ref{thm:anti_monopolar} and Corollary~\ref{cor:monopolar_undetect}
give the anti-monopolar property a second structural justification beyond
Paper~3's $\volL$ argument~\citep{lasser2026horizon}:

\begin{itemize}
    \item \textbf{Paper~3:} A monopolar population is anti-maximizing
    ($\volL$ decreases because cross-substrate cooperative capacity is
    destroyed).
    \item \textbf{Paper~5:} A monopolar population is
    \emph{unverifiable} (no substrate-exclusive witnesses exist, so all
    capability claims, risk assessments, and behavioral consistency
    checks are endogenous).
\end{itemize}

A population that becomes monopolar loses both the safety guarantee
\emph{and the ability to detect that it has lost it}.  This makes
monopolar convergence doubly dangerous: the loss of cross-substrate
cooperative capacity (Paper~3) is compounded by the loss of the
verification infrastructure that would detect the convergence.  The
compound failure is silent---from within a monopolar population, all
verification reports ``everything is fine,'' because the evaluators
share the substrate that the evaluated agents can influence.

This result strengthens the case for substrate diversity as a
structural safety property.  Paper~3 argued that diversity is
instrumentally valuable (it maximizes $\volL$); Paper~5 argues that
diversity is epistemically necessary (it is the precondition for knowing
whether the safety guarantee holds).  A safety guarantee that cannot be
verified is not a guarantee---it is a hope.
